\chapter{Les fondements de la géométrie plane}
\label{corriges:12}
\section*{Exercice 12.1 -- Angles d'un triangle}
Dans l'isocèle, $\widehat B=\widehat C$. Leur somme vaut $180-44=136\degre$, donc chacun vaut $68\degre$. Dans l'autre triangle, le troisième angle vaut $180-35-55=90\degre$ : il est rectangle. Les deux conclusions utilisent la somme des angles d'un triangle.

\section*{Exercice 12.2 -- Construire et décider si c'est possible}
Pour $5,7,10$, $5+7>10$ : tracer un segment de $10\unite{cm}$, puis des cercles de rayons $5$ et $7\unite{cm}$ centrés à ses extrémités. Une intersection est le troisième sommet ; l'autre donne le triangle retourné. Pour $3,4,8$, $3+4=7<8$ contredit l'inégalité triangulaire : construction impossible.

\section*{Exercice 12.3 -- Médiatrice et cercle}
Tracer deux arcs de même rayon supérieur à $AB/2$, centrés en $A$ et $B$, puis relier leurs intersections. Elles sont équidistantes de $A$ et $B$, donc appartiennent à la médiatrice.
Pour $M$, les égalités donnent $MA=MB=MC$. Le cercle de centre $M$ et de rayon $MA$ passe par les trois points : $M$ est le centre de leur cercle circonscrit, donc l'intersection des médiatrices.

\section*{Exercice 12.4 -- Quel quadrilatère ?}
Le même milieu établit d'abord que le quadrilatère est un parallélogramme. Ses diagonales égales le caractérisent alors comme un rectangle. Si elles sont aussi perpendiculaires, il est également un losange : c'est donc un carré. Sans l'hypothèse du milieu commun, les autres données ne suffiraient pas seules.
